\section{模型评价、改进与推广}

\subsection{模型优点}

本文模型的第一个优点是数据到模型的映射清晰。地块、作物、季次和年份被统一为可行三元组集合 $\Omega$，附件中的地块规则不再以文字形式散落在模型外，而是直接转化为变量是否存在。这种处理减少了无效变量，也使后续约束更容易检查。

第二，模型能够同时表达农业制度和经济目标。面积上限、最小种植比例、连作禁止、豆类轮作和作物集中度约束保证方案不只追求账面收益，也满足基本农艺要求。对于数学建模而言，这一点比单纯最大化利润更重要，因为农业方案最终必须能落到地块和季次上执行。

第三，问题二和问题三逐步加入现实因素。随机规划处理需求、价格、成本和亩产波动，CVaR刻画尾部风险，作物相关性模型进一步考虑替代和互补机制。三问之间不是孤立求解，而是从确定性基准到随机风险，再到作物关联逐层扩展，便于解释结果变化的来源。

\subsection{模型不足}

本文的主要不足在于问题二和问题三采用了拓扑固定法。该方法使大规模随机规划可以在普通计算环境下求解，并能获得固定拓扑下的全局最优解，但它没有重新优化全部0-1种植指示变量。因此，问题二和问题三的结果应理解为在问题一半价拓扑基础上的风险调整与结构扩展方案，而不是全空间混合整数随机规划的全局最优方案。

第二，作物替代关系主要由统计相关性近似识别。Spearman相关系数能够反映市场指标的同步或反向变化，但不能完全代表真实消费替代弹性。由于题目未给出消费者选择、市场价格联动或外部需求曲线，本文采用相关性和聚类作为可解释的近似方法。

第三，灾害模型只考虑露天耕地的随机减产，尚未细分灾害类型、地块空间相邻性和灾害持续时间。该处理符合题目给定随机规则，但在更真实的农业风险管理中仍可进一步细化。

\subsection{改进方向}

若计算资源允许，可以在问题二中对部分关键作物或关键地块重新开放0-1变量，形成“局部拓扑可变”的随机混合整数模型。这样既不会像全量MILP那样规模过大，又能检验当前固定拓扑是否限制了随机环境下的收益提升空间。

在作物关系方面，可引入历史市场价格、销量或替代消费数据，估计更接近经济意义的需求弹性矩阵。若能获得豆类轮作对后茬产量和成本的实测影响，也可以把互补关系从固定折减系数扩展为作物和地块类型相关的参数。

在风险分析方面，可进一步区分干旱、暴雨和病虫害等灾害类型，并建立地块空间相关的灾害情景。这样能够更准确地评估露天耕地集中种植时可能出现的系统性风险。

\subsection{模型推广}

本文模型可推广到其他多地块、多作物、多季次的农业规划场景。若用于其他地区，只需替换地块集合、作物集合、适宜种植规则、历史统计数据和市场不确定性分布。若用于设施农业或订单农业，还可以把大棚容量、合同销售量、最低供应量和劳动力约束加入现有框架。模型结构保持不变，主要变化集中在参数表和约束集合上。
